% lit_execution.tex -- Literature 2.5: the execution layer and force estimation (drafted 12 Oct 2026 from the
% thesis plan of 11 Oct; the steady-error paragraph is from control.tex). Entries the plan lists that are not yet in
% the bibliography were added 12 Oct (Tagliabue 2019, de Marina and Smeur 2019, Bulka 2022).

\subsection{Execution layer and force estimation}
\label{sec:lit_execution}

An execution layer sits between a planner's references and each vehicle's actuation, and in cooperative transport
it has one job the single-vehicle case does not: the link's pull. In Sun et al.\ \cite{tudelft} an onboard INDI
controller \cite{smeur2016indi} turns the planned trajectories into motor commands and compensates the cable force
from the vehicle's own accelerometer, without a model of the cable: centralised planning paired with robust
local execution, the fast dynamics handled onboard and the slow coordination centrally. Estimating the link's
force rather than measuring it is the other route: an adaptive estimate of an unknown load in a decentralised
passivity-based scheme \cite{klausen2020cooperative}, or a model-based observer of the external wrench
\cite{tagliabue2019robust}; de Marina and Smeur \cite{demarina2019flexible} and Bulka et al.\ \cite{bulka2022experiments} fly a shared load with the execution onboard each rotorcraft. Each of these needs code on the vehicle.

The vehicles of this thesis run unmodified Betaflight, a hobby flight controller that closes a body-rate loop and
accepts a collective throttle and three rate commands over a radio link, as a pilot would send them. Nothing can
be added onboard, so the execution layer is a tracking MPC per vehicle on the ground, with the rod's pull entering
its model as a known input along the horizon: from the planner's own tensions, or measured from the vehicle's
motion-capture acceleration minus its modelled thrust (Section~\ref{sec:tracker}). What the measured pull buys
over the planned one on this hardware is the first research question (Section~\ref{sec:rq}).

\paragraph{Steady-error rejection.} Without an integrating element, a constant disturbance or a model error
leaves a steady error in what the planner controls. Offset-free MPC removes it by adding a constant disturbance
to the prediction model and estimating it from the measurements \cite{muske2002disturbance,
pannocchia2003disturbance, maeder2009linear}; by the internal model principle \cite{francis1976internal} a
constant disturbance is then rejected with no steady error. The disturbance can be modelled at the input or at
the output, which changes the transient and not the steady state \cite{pannocchia2007combined}.
Section~\ref{sec:disturbance} builds the planner's integral action on this, as a force on the payload in the
planner's own model, and Appendix~\ref{sec:app_methods} records the estimators that were built and flown against
it.
